\scanpage{285}
\mspage{201}
% archive sheet number: 284
\uline{Corollaire} 4.8.5. Soit $C$ une \deleted{petite} catégorie équivalente à une petite catégorie, \deleted{$\unsure{M}$} $\unsure{\pi|\kappa|k}_0$ un cardinal infini, $M=\mathrm{Ind}_{\pi_0}(C)\subset C^{\wedge}$. Alors

$1^\circ)$ $\pi_0$ est un cardinal \deleted{\uline{régulier}}\inserted{utile, et $=$ \uline{régulier}} pour la \deleted{\unsure{filtr}} feg canonique $(M_\pi)_{\pi\geq\pi_0}$ de $M$.

$2^\circ)$ $M_{\pi_0}\unsure{=|\simeq}C^{\unsure{\natural|4}}$ (sous-catégorie pleine de \inserted{$M=$}$\mathrm{Ind}_{\pi_0}C$ formée des facteurs directs dans $M$ des objets de $C$), et $M_{\pi_0}$ est une enveloppe de \unsure{Karoubi} de $C$. \unsure{a}

% slanted
\margin{\uline{4.8.5.2} \unsure{Expliciter|Explicitons} aussi ceci : Soit $u:C\longrightarrow C'$ un foncteur, \unsure{avec|où} $C,C'$ \unsure{deux|sont} petites catégories. Conditions équivalentes\unsure{:} \par a) $u$ induit $\mathrm{Kar}\,C\xrightarrow{\approx}\mathrm{Kar}\,C'$. \par a') $u$ pl. fid., et tout objet de $C'$ \unsure{isom.} à un facteur \unsure{direct|dir.} d'un \unsure{objet} de l'image essentielle\unsure{.} \par b) $\mathrm{Ind}(u):\mathrm{Ind}(C)\longrightarrow\mathrm{Ind}(C')$ est une équiv. \par c) $\mathrm{Ind}_{\pi_0}(u):\mathrm{Ind}_{\pi_0}(C)\longrightarrow\mathrm{Ind}_{\pi_0}(C')$ est une équivalence ($\pi_0$ un cardinal infini donné) \par d) $u_!:C^{\wedge}\longrightarrow C'^{\wedge}$ est une équivalence \par d') $u_*:C'^{\wedge}\longrightarrow C^{\wedge}$ est une équivalence \par d'') $u^{\unsure{\wedge|*}}:C'^{\wedge}\longrightarrow C^{\wedge}$ est une équivalence}

\margin{\uline{Corollaire} 4.8.5.1. Si $C$ est Karoubienne, alors $M_{\pi_0}\subset C^{\wedge}$ est la catégorie formée des foncteurs \uline{représentables}, donc $C\xrightarrow{\approx}M_{\pi_0}$.}

\uline{Dém} \inserted{$1^\circ)$} On sait déjà que

\[ \mathrm{Ind}_{\pi_0}M_{\pi_0}\longrightarrow M \]

est pleinement fidèle, et pour voir que c'est ess. surjectif, on note que l'on a

\[ M=\mathrm{Ind}_{\pi_0}C\mathrel{\unsure{\leftrightarrow|\hookrightarrow}}\mathrm{Ind}_{\pi_0}M_{\unsure{0|\pi_0}}\longrightarrow M, \]

le composé étant isom. \deleted{\unsure{donc}} au foncteur identique, est ess. surjectif, donc aussi \unsure{cas.|ceci.} Cela prouve donc que $\pi_0$ est un cardinal utile pour la filtration envisagée, prouvons que ce cardinal est $=$ régulier, donc que pour

\[ F\in\mathrm{Ob}\,M_{\pi},\qquad\text{avec }\left[\begin{aligned}
\pi\ \deleted{\unsure{\geq}}&>\pi_0\\
\pi^{\pi_0}&\geq\pi_1,
\end{aligned}\right. \]

on a

\[ F=\varinjlim_I \deleted{\unsure{X}}x_i \quad \left|\begin{aligned}
&x_i\in\mathrm{Ob}\,M_{\pi_0}\\
&I\ \text{grd devant }\pi_0\\
&\mathrm{Card}\, I\leq\pi^{\pi_0}.
\end{aligned}\right. \]

\deleted{\unsure{pour} $\pi^{\pi_0}\geq\pi_1$\unsure{.}}
